% \documentclass[a4paper,13pt,3p,twoside]{report}
% \usepackage{scrextend}
% \changefontsizes{13pt}
% \usepackage[utf8]{vietnam}
% \usepackage[top=2cm, bottom=2cm, left=3.5cm, right=2.5cm]{geometry}
% \usepackage{xurl}
% \usepackage{appendix}
% \usepackage{babel}
% \usepackage{xcolor}
% \usepackage{outlines}
% \usepackage{graphicx} % Cho phép chèn hỉnh ảnh
% \usepackage{fancybox} % Tạo khung box
% \usepackage{indentfirst} % Thụt đầu dòng ở dòng đầu tiên trong đoạn
% \usepackage{amsthm} % Cho phép thêm các môi trường định nghĩa
% \usepackage{latexsym} % Các kí hiệu toán học
% \usepackage{amsmath} % Hỗ trợ một số biểu thức toán học
% \usepackage{amssymb} % Bổ sung thêm kí hiệu về toán học
% \usepackage{amsbsy} % Hỗ trợ các kí hiệu in đậm
% \usepackage{times} % Chọn font Time New Romans
% \usepackage{array} % Tạo bảng array
% \usepackage{enumitem} % Cho phép thay đổi kí hiệu của list
% \usepackage{subfiles} % Chèn các file nhỏ, giúp chia các chapter ra nhiều file hơn
% \usepackage{enumitem}
% \newlist{romanlist}{enumerate}{1}
% \setlist[romanlist]{label=(\roman*), leftmargin=2.2em, itemsep=2pt}
% \usepackage{titlesec} % Giúp chỉnh sửa các tiêu đề, đề mục như chương, phần,..
% \usepackage{titletoc}

% % Define the designspec environment
% \newenvironment{designspec}[1]{%
%   \begin{center}
%     \begin{tabular}{|p{0.95\textwidth}|}
%     \hline
%     \textbf{#1} \\
%     \hline
%     \begin{minipage}{0.93\textwidth}
% }{%
%     \end{minipage} \\
%     \hline
%     \end{tabular}
%   \end{center}
% }
% \usepackage{chngcntr} % Dùng để thiết lập lại cách đánh số caption,..
% \usepackage{pdflscape} % Đưa các bảng có kích thước đặt theo chiều ngang giấy
% \usepackage{afterpage}
% \usepackage[ruled,vlined]{algorithm2e}  % Hỗ trợ viết các giải thuật
% \usepackage{capt-of} % Cho phép sử dụng caption lớn đối với landscape page
% \usepackage{multirow} % Merge cells
% \usepackage{fancyhdr} % Cho phép tùy biến header và footer
% % \usepackage[natbib,backend=biber,style=ieee]{biblatex} % Giúp chèn tài liệu tham khảo

% \usepackage[font=small,labelfont=bf]{caption}

% \usepackage{listings}
% \usepackage{float}
% \usepackage{subcaption}

% \usepackage[nonumberlist, nopostdot, nogroupskip, acronym]{glossaries}
% % \makeglossaries
% \usepackage{glossary-superragged}
% \setlength{\glsdescwidth}{10.5cm}
% \setlength{\glspagelistwidth}{0pt}
% \setglossarystyle{superraggedheaderborder}
% \usepackage{setspace}
% \usepackage{parskip}

% % package content table
% \usepackage{tocbasic}

% \usepackage{blindtext}


% % ===================================================

% % \renewcommand{\bibname}{Danh_sach_tai_lieu_tham_khao} 
% \usepackage[backend=bibtex,style=ieee]{biblatex}  %backend=biber is 'better'

% \addbibresource{reference.bib} % chèn file chứa danh mục tài liệu tham khảo vào 

% \include{lstlisting} % Phần này cho phép chèn code và formatting code như C, C++, Python
% \makenoidxglossaries 
% % %\makeglossaries
% \makenoidxglossaries

% % Danh mục thuật ngữ và từ viết tắt
% \newglossaryentry{iaas}{
%     type=\acronymtype,
%     name={IaaS},
%     description={Dịch vụ hạ tầng},
%     first={Dịch vụ hạ tầng (Infrastructure  as  a  Service - IaaS)}
% }
% \newglossaryentry{API}{
%     type=\acronymtype,
%     name={API},
%     description={Giao diện lập trình ứng dụng (Application Programming Interface)},
%     first={API}
% }
% \newglossaryentry{EUD}{
%     type=\acronymtype,
%     name={EUD},
%     description={Phát triển ứng dụng người dùng cuối(End-User Development)},
%     first={End-User Development}
% }
% \newglossaryentry{GWT}{
%     type=\acronymtype,
%     name={GWT},
%     description={Công cụ lập trình Javascript bằng Java của Google (Google Web Toolkit)},
%     first={Công cụ lập trình Javascript bằng Java của Google (Google Web Toolkit}
% }
% \newglossaryentry{HTML}{
%     type=\acronymtype,
%     name={HTML},
%     description={Ngôn ngữ đánh dấu siêu văn bản (HyperText Markup Language)},
%     first={Dịch vụ hạ tầng (Infrastructure  as  a  Service - IaaS)}
% }

% ---- Core system concepts -------------------------------------------
\newacronym{ids}{IDS}{Intrusion Detection System}
\newacronym{ips}{IPS}{Intrusion Prevention System}
\newacronym{waf}{WAF}{Web Application Firewall}
\newacronym{siem}{SIEM}{Security Information and Event Management}

% ---- Machine-learning models ----------------------------------------
\newacronym{ml}{ML}{Machine Learning}
\newacronym{rf}{RF}{Random Forest}
\newacronym{lr}{LR}{Logistic Regression}
\newacronym{lof}{LOF}{Local Outlier Factor}
\newacronym{if}{IF}{Isolation Forest}
\newacronym{ocsvm}{OCSVM}{One-Class Support Vector Machine}

% ---- Web protocols and formats --------------------------------------
\newacronym{http}{HTTP}{HyperText Transfer Protocol}
\newacronym{https}{HTTPS}{HyperText Transfer Protocol Secure}
\newacronym{url}{URL}{Uniform Resource Locator}
\newacronym{uri}{URI}{Uniform Resource Identifier}
\newacronym{html}{HTML}{HyperText Markup Language}
\newacronym{json}{JSON}{JavaScript Object Notation}
\newacronym{csv}{CSV}{Comma-Separated Values}
\newacronym{api}{API}{Application Programming Interface}

% ---- Attack types ---------------------------------------------------
\newacronym{sqli}{SQLi}{SQL Injection}
\newacronym{sql}{SQL}{Structured Query Language}
\newacronym{xss}{XSS}{Cross-Site Scripting}
\newacronym{lfi}{LFI}{Local File Inclusion}
\newacronym{rfi}{RFI}{Remote File Inclusion}
\newacronym{owasp}{OWASP}{Open Worldwide Application Security Project}

% ---- Datasets -------------------------------------------------------
\newacronym{csic}{CSIC}{Spanish Research Council HTTP Dataset (CSIC 2010)}
\newacronym{ecml}{ECML/PKDD}{European Conference on Machine Learning and
  Principles and Practice of Knowledge Discovery in Databases}

% ---- Evaluation metrics ---------------------------------------------
\newacronym{tp}{TP}{True Positive}
\newacronym{fp}{FP}{False Positive}
\newacronym{tn}{TN}{True Negative}
\newacronym{fn}{FN}{False Negative}
\newacronym{tpr}{TPR}{True Positive Rate (Recall)}
\newacronym{fpr}{FPR}{False Positive Rate}
\newacronym{roc}{ROC}{Receiver Operating Characteristic}
\newacronym{auc}{AUC}{Area Under the Curve}

% ---- Deployment / infrastructure ------------------------------------
\newacronym{keda}{KEDA}{Kubernetes Event-Driven Autoscaling}
\newacronym{elk}{ELK}{Elasticsearch, Logstash and Kibana}
\newacronym{cms}{CMS}{Content Management System}


% % ===================================================


% \fancypagestyle{plain}{%
% \fancyhf{} % clear all header and footer fields
% \fancyfoot[RO,RE]{\thepage} %RO=right odd, RE=right even
% \renewcommand{\headrulewidth}{0pt}
% \renewcommand{\footrulewidth}{0pt}}

% \setlength{\headheight}{10pt}

% \def \TITLE{GRADUATION THESIS}
% \def \AUTHOR{Trần Văn A}

% % ===================================================
% \titleformat{\chapter}[hang]{\centering\bfseries}{CHAPTER \thechapter.\ }{0pt}{}[]

% \titleformat 
%     {\chapter} % command
%     [hang] % shape
%     {\centering\bfseries} % format
%     {CHAPTER \thechapter.\ } % label
%     {0pt} %sep
%     {} % before
%     [] % after
% \titlespacing*{\chapter}{0pt}{-20pt}{20pt}

% \titleformat
%     {\section} % command
%     [hang] % shape
%     {\bfseries} % format
%     {\thechapter.\arabic{section}\ \ \ \ } % label
%     {0pt} %sep
%     {} % before
%     [] % after
% \titlespacing{\section}{0pt}{\parskip}{0.5\parskip}

% \titleformat
%     {\subsection} % command
%     [hang] % shape
%     {\bfseries} % format
%     {\thechapter.\arabic{section}.\arabic{subsection}\ \ \ \ } % label
%     {0pt} %sep
%     {} % before
%     [] % after
% \titlespacing{\subsection}{30pt}{\parskip}{0.5\parskip}

% \renewcommand\thesubsubsection{\alph{subsubsection}}
% \titleformat
%     {\subsubsection} % command
%     [hang] % shape
%     {\bfseries} % format
%     {\alph{subsubsection}, \ } % label
%     {0pt} %sep
%     {} % before
%     [] % after
% \titlespacing{\subsubsection}{50pt}{\parskip}{0.5\parskip}

% % \newcommand{\titlesize}{\fontsize{18pt}{23pt}\selectfont}
% % \newcommand{\subtitlesize}{\fontsize{16pt}{21pt}\selectfont}
% % \titleclass{\part}{top}
% % \titleformat{\part}[display]
% %   {\normalfont\huge\bfseries}{\centering}{20pt}{\Huge\centering}
% % \titlespacing{\part}{0pt}{em}{1em}
% % \titlespacing{\section}{0pt}{\parskip}{0.5\parskip}
% % \titlespacing{\subsection}{0pt}{\parskip}{0.5\parskip}
% % \titlespacing{\subsubsection}{0pt}{\parskip}{0.5\parskip}



% % ===================================================
% \usepackage{hyperref}
% \hypersetup{pdfborder = {0 0 0}}
% \hypersetup{pdftitle={\TITLE},
% 	pdfauthor={\AUTHOR}}
	
% \usepackage[all]{hypcap} % Cho phép tham chiếu chính xác đến hình ảnh và bảng biểu

% \graphicspath{{figures/}{../figures/}} % Thư mục chứa các hình ảnh

% \counterwithin{figure}{chapter} % Đánh số hình ảnh kèm theo chapter. Ví dụ: Hình 1.1, 1.2,..

% \title{\bf \TITLE}
% \author{\AUTHOR}

% \setcounter{secnumdepth}{3} % Cho phép subsubsection trong report
% % \setcounter{tocdepth}{3} % Chèn subsubsection vào bảng mục lục

% \theoremstyle{definition}
% \newtheorem{example}{Example}[chapter] % Định nghĩa môi trường ví dụ

% \onehalfspacing
% \setlength{\parskip}{6pt}
% \setlength{\parindent}{15pt}



% % =========================== BODY ===============
% \begin{document}
% % \newgeometry{top=2cm, bottom=2cm, left=2cm, right=2cm}
% % \subfile{Cover} % Phần bìa
% \subfile{Cover} % Phần bìa lót
% % \restoregeometry

% % ===================================================
% \pagestyle{empty} % Header và footer rỗng
% %\newpage
% %\pagenumbering{gobble} % Xóa page numbering ở cuối trang
% %\subfile{chapters/0_1_subject.tex}

% % \pagestyle{empty} % Header và footer rỗng
% \newpage
% \pagenumbering{gobble} % Xóa page numbering ở cuối trang
% \subfile{Chapter/0_2_Acknowledgement}

% % \pagestyle{empty} % Header và footer rỗng
% \newpage
% \pagenumbering{gobble} % Xóa page numbering ở cuối trang
% \subfile{Chapter/0_3_Abstract.tex}


% % ===================================================
% % \pagestyle{empty} % Header và footer rỗng
% \newpage
% \pagenumbering{gobble} % Xóa page numbering ở cuối trang
% \renewcommand*\contentsname{TABLE OF CONTENTS}
% \titlecontents{chapter}
%     [0.0cm]             % left margin
%     {\bfseries\vspace{0.3cm}}                  % above code
%     {{\bfseries{\scshape} CHAPTER \thecontentslabel.\ }} % numbered format
%     {}         % unnumbered format
%     {\titlerule*[0.3pc]{.}\contentspage}         % filler-page-format, e.g dots
    
% \titlecontents{section}
%     [0.0cm]             % left margin
%     {\vspace{0.3cm}}                  % above code
%     {\thecontentslabel \ } % numbered format
%     {}         % unnumbered format
%     {\titlerule*[0.3pc]{.}\contentspage}         % filler-page-format, e.g dots
    
% \titlecontents{subsection}
%     [1.0cm]             % left margin
%     {\vspace{0.3cm}}                  % above code
%     {\thecontentslabel \ } % numbered format
%     {}         % unnumbered format
%     {\titlerule*[0.3pc]{.}\contentspage}         % filler-page-format, e.g dots

%  % Tạo mục lục tự động
% % \addtocontents{toc}{\protect\thispagestyle{empty}}
% % \tableofcontents 
% % \thispagestyle{empty}
% % \cleardoublepage

% % \pagenumbering{roman}
% % %Tạo danh mục hình vẽ.
% % \renewcommand{\listfigurename}{LIST OF FIGURES}
% % {\let\oldnumberline\numberline
% % \renewcommand{\numberline}{Figure~\oldnumberline}
% % \listoffigures} 
% % % \phantomsection\addcontentsline{toc}{section}{\numberline {} DANH MỤC HÌNH VẼ}
% % \newpage


% %  %Tạo danh mục bảng biểu.
% % \renewcommand{\listtablename}{LIST OF TABLES}
% % {\let\oldnumberline\numberline
% % \renewcommand{\numberline}{Table~\oldnumberline}
% % \listoftables}
% % % \phantomsection\addcontentsline{toc}{section}{\numberline {} DANH MỤC BẢNG BIỂU}

% % % \renewcommand*{\glossaryname}{Danh sách thuật ngữ}
% % \renewcommand*{\acronymname}{LIST OF ABBREVIATIONS}
% % \renewcommand*{\entryname}{Abbreviation}
% % \renewcommand*{\descriptionname}{Definition}
% % \glsaddall
% % \printnoidxglossary[type=\acronymtype, title={LIST OF ABBREVIATIONS}, nonumberlist]
% % \phantomsection\addcontentsline{toc}{chapter}{LIST OF ABBREVIATIONS}
% % \newpage
% % % \phantomsection\addcontentsline{toc}{section}{\numberline {} DANH MỤC THUẬT NGỮ VÀ TỪ VIẾT TẮT}

% % \renewcommand\appendixname{APPENDIX}
% % \renewcommand\appendixpagename{APPENDIX}
% % \renewcommand\appendixtocname{APPENDIX}

% % \renewcommand{\figurename}{Figure}
% % \renewcommand{\tablename}{Table}
% % \renewcommand{\chaptername}{CHAPTER}

% % % ===================================================


% % \newpage
% % \pagenumbering{arabic}

% % \pagestyle{fancy}
% % \fancyhf{}
% % \fancyhead[RE, LO]{\leftmark}
% % %\fancyhead[LE]{\rightmark}
% % \fancyfoot[RE, LO]{\thepage}

% \addtocontents{toc}{\protect\thispagestyle{empty}}
% \tableofcontents
% \thispagestyle{empty}
% \clearpage                       % <-- ĐỔI từ \cleardoublepage

% \pagenumbering{roman}

% % ---- LIST OF FIGURES ----
% \renewcommand{\listfigurename}{LIST OF FIGURES}
% {\let\oldnumberline\numberline
%  \renewcommand{\numberline}{Figure~\oldnumberline}
%  \listoffigures}
% \clearpage                       % <-- ĐỔI từ \newpage

% % ---- LIST OF TABLES ----
% \renewcommand{\listtablename}{LIST OF TABLES}
% {\let\oldnumberline\numberline
%  \renewcommand{\numberline}{Table~\oldnumberline}
%  \listoftables}
% \clearpage                       % <-- THÊM

% % ---- LIST OF ABBREVIATIONS ----
% \renewcommand*{\acronymname}{LIST OF ABBREVIATIONS}
% \renewcommand*{\entryname}{Abbreviation}
% \renewcommand*{\descriptionname}{Definition}
% \glsaddall                       % <-- chỉ MỘT lần (bỏ cái lặp)
% \printnoidxglossary[type=\acronymtype, title={LIST OF ABBREVIATIONS}, nonumberlist]
% \phantomsection\addcontentsline{toc}{chapter}{LIST OF ABBREVIATIONS}
% \clearpage

% \renewcommand\appendixname{APPENDIX}
% \renewcommand\appendixpagename{APPENDIX}
% \renewcommand\appendixtocname{APPENDIX}
% \renewcommand{\figurename}{Figure}
% \renewcommand{\tablename}{Table}
% \renewcommand{\chaptername}{CHAPTER}

% \pagenumbering{arabic}           % <-- BỎ dòng \newpage ngay trước nó
% \pagestyle{fancy}
% \fancyhf{}
% \fancyhead[RE, LO]{\leftmark}
% \fancyfoot[RE, LO]{\thepage}

% \chapter{INTRODUCTION}
% \subfile{Chapter/1_Introduction} % Phần mở đầu

% \newpage
% %\pagestyle{fancy} % Áp dụng header và footer
% \chapter{BACKGROUND AND RELATED WORK}
% \subfile{Chapter/2_Literature_review}


% \newpage
% %\pagestyle{fancy} % Áp dụng header và footer
% \chapter{PROPOSED SOLUTION}
% \subfile{Chapter/3_Methodology}

% \newpage
% %\pagestyle{fancy} % Áp dụng header và footer
% \chapter{EVALUATION AND RESULTS}
% \subfile{Chapter/5_Numerical_results}

% \newpage
% \chapter{CONCLUSIONS AND FUTURE WORK}
% \subfile{Chapter/6_Conclusions}

% % ===================================================
% \newpage
% \renewcommand\bibname{REFERENCE}
% \printbibliography
% \phantomsection\addcontentsline{toc}{chapter}{REFERENCE}

% \end{document}

\documentclass[a4paper,13pt,3p,twoside]{report}
\usepackage{scrextend}
\changefontsizes{13pt}
\usepackage[utf8]{vietnam}
\usepackage[top=2cm, bottom=2cm, left=3.5cm, right=2.5cm]{geometry}
\usepackage{xurl}
\usepackage{appendix}
\usepackage{babel}
\usepackage{xcolor}
\usepackage{outlines}
\usepackage{graphicx} % Cho phép chèn hỉnh ảnh
\usepackage{fancybox} % Tạo khung box
\usepackage{indentfirst} % Thụt đầu dòng ở dòng đầu tiên trong đoạn
\usepackage{amsthm} % Cho phép thêm các môi trường định nghĩa
\usepackage{latexsym} % Các kí hiệu toán học
\usepackage{amsmath} % Hỗ trợ một số biểu thức toán học
\usepackage{amssymb} % Bổ sung thêm kí hiệu về toán học
\usepackage{amsbsy} % Hỗ trợ các kí hiệu in đậm
\usepackage{times} % Chọn font Time New Romans
\usepackage{array} % Tạo bảng array
\usepackage{enumitem} % Cho phép thay đổi kí hiệu của list
\usepackage{subfiles} % Chèn các file nhỏ, giúp chia các chapter ra nhiều file hơn
\usepackage{enumitem}
\newlist{romanlist}{enumerate}{1}
\setlist[romanlist]{label=(\roman*), leftmargin=2.2em, itemsep=2pt}
\usepackage{titlesec} % Giúp chỉnh sửa các tiêu đề, đề mục như chương, phần,..
\usepackage{titletoc}

% Define the designspec environment
\newenvironment{designspec}[1]{%
  \begin{center}
    \begin{tabular}{|p{0.95\textwidth}|}
    \hline
    \textbf{#1} \\
    \hline
    \begin{minipage}{0.93\textwidth}
}{%
    \end{minipage} \\
    \hline
    \end{tabular}
  \end{center}
}
\usepackage{chngcntr} % Dùng để thiết lập lại cách đánh số caption,..
\usepackage{pdflscape} % Đưa các bảng có kích thước đặt theo chiều ngang giấy
\usepackage{afterpage}
\usepackage[ruled,vlined]{algorithm2e}  % Hỗ trợ viết các giải thuật
\usepackage{capt-of} % Cho phép sử dụng caption lớn đối với landscape page
\usepackage{multirow} % Merge cells
\usepackage{fancyhdr} % Cho phép tùy biến header và footer
% \usepackage[natbib,backend=biber,style=ieee]{biblatex} % Giúp chèn tài liệu tham khảo

\usepackage[font=small,labelfont=bf]{caption}

\usepackage{listings}
\usepackage{float}
\usepackage{subcaption}

% ===================== GLOSSARY / ABBREVIATIONS =====================
\usepackage[nonumberlist, nopostdot, nogroupskip, acronym]{glossaries}
% \makeglossaries   % KHONG dung: ta in bang \printnoidxglossary
\usepackage{glossary-longbooktabs} % style longtable + booktabs (tu nap longtable & booktabs)
\setlength{\glsdescwidth}{11cm}    % do rong cot "Definition" (vua trong le ~15cm)
\setglossarystyle{long-booktabs}   % bang sach, tu ngat trang, khong tran
% Ep tieu de glossary thanh MOT dong can giua (giong LIST OF FIGURES/TABLES):
% khong nhay trang, khong tao khoang trang 50pt phia tren tieu de.
\renewcommand*{\glossarysection}[2][]{%
  \vspace*{-10pt}{\centering\bfseries #2\par}\vspace{18pt}}
% ===================================================================
\usepackage{setspace}
\usepackage{parskip}

% package content table
\usepackage{tocbasic}

\usepackage{blindtext}


% ===================================================

% \renewcommand{\bibname}{Danh_sach_tai_lieu_tham_khao}
\usepackage[backend=bibtex,style=ieee]{biblatex}  %backend=biber is 'better'

\addbibresource{reference.bib} % chèn file chứa danh mục tài liệu tham khảo vào

\include{lstlisting} % Phần này cho phép chèn code và formatting code như C, C++, Python
\makenoidxglossaries
% %\makeglossaries
% \makenoidxglossaries

% % Danh mục thuật ngữ và từ viết tắt
% \newglossaryentry{iaas}{
%     type=\acronymtype,
%     name={IaaS},
%     description={Dịch vụ hạ tầng},
%     first={Dịch vụ hạ tầng (Infrastructure  as  a  Service - IaaS)}
% }
% \newglossaryentry{API}{
%     type=\acronymtype,
%     name={API},
%     description={Giao diện lập trình ứng dụng (Application Programming Interface)},
%     first={API}
% }
% \newglossaryentry{EUD}{
%     type=\acronymtype,
%     name={EUD},
%     description={Phát triển ứng dụng người dùng cuối(End-User Development)},
%     first={End-User Development}
% }
% \newglossaryentry{GWT}{
%     type=\acronymtype,
%     name={GWT},
%     description={Công cụ lập trình Javascript bằng Java của Google (Google Web Toolkit)},
%     first={Công cụ lập trình Javascript bằng Java của Google (Google Web Toolkit}
% }
% \newglossaryentry{HTML}{
%     type=\acronymtype,
%     name={HTML},
%     description={Ngôn ngữ đánh dấu siêu văn bản (HyperText Markup Language)},
%     first={Dịch vụ hạ tầng (Infrastructure  as  a  Service - IaaS)}
% }

% ---- Core system concepts -------------------------------------------
\newacronym{ids}{IDS}{Intrusion Detection System}
\newacronym{ips}{IPS}{Intrusion Prevention System}
\newacronym{waf}{WAF}{Web Application Firewall}
\newacronym{siem}{SIEM}{Security Information and Event Management}

% ---- Machine-learning models ----------------------------------------
\newacronym{ml}{ML}{Machine Learning}
\newacronym{rf}{RF}{Random Forest}
\newacronym{lr}{LR}{Logistic Regression}
\newacronym{lof}{LOF}{Local Outlier Factor}
\newacronym{if}{IF}{Isolation Forest}
\newacronym{ocsvm}{OCSVM}{One-Class Support Vector Machine}

% ---- Web protocols and formats --------------------------------------
\newacronym{http}{HTTP}{HyperText Transfer Protocol}
\newacronym{https}{HTTPS}{HyperText Transfer Protocol Secure}
\newacronym{url}{URL}{Uniform Resource Locator}
\newacronym{uri}{URI}{Uniform Resource Identifier}
\newacronym{html}{HTML}{HyperText Markup Language}
\newacronym{json}{JSON}{JavaScript Object Notation}
\newacronym{csv}{CSV}{Comma-Separated Values}
\newacronym{api}{API}{Application Programming Interface}

% ---- Attack types ---------------------------------------------------
\newacronym{sqli}{SQLi}{SQL Injection}
\newacronym{sql}{SQL}{Structured Query Language}
\newacronym{xss}{XSS}{Cross-Site Scripting}
\newacronym{lfi}{LFI}{Local File Inclusion}
\newacronym{rfi}{RFI}{Remote File Inclusion}
\newacronym{owasp}{OWASP}{Open Worldwide Application Security Project}

% ---- Datasets -------------------------------------------------------
\newacronym{csic}{CSIC}{Spanish Research Council HTTP Dataset (CSIC 2010)}
\newacronym{ecml}{ECML/PKDD}{European Conference on Machine Learning and
  Principles and Practice of Knowledge Discovery in Databases}

% ---- Evaluation metrics ---------------------------------------------
\newacronym{tp}{TP}{True Positive}
\newacronym{fp}{FP}{False Positive}
\newacronym{tn}{TN}{True Negative}
\newacronym{fn}{FN}{False Negative}
\newacronym{tpr}{TPR}{True Positive Rate (Recall)}
\newacronym{fpr}{FPR}{False Positive Rate}
\newacronym{roc}{ROC}{Receiver Operating Characteristic}
\newacronym{auc}{AUC}{Area Under the Curve}

% ---- Deployment / infrastructure ------------------------------------
\newacronym{keda}{KEDA}{Kubernetes Event-Driven Autoscaling}
\newacronym{elk}{ELK}{Elasticsearch, Logstash and Kibana}
\newacronym{cms}{CMS}{Content Management System}


% ===================================================


\fancypagestyle{plain}{%
\fancyhf{} % clear all header and footer fields
\fancyfoot[RO,RE]{\thepage} %RO=right odd, RE=right even
\renewcommand{\headrulewidth}{0pt}
\renewcommand{\footrulewidth}{0pt}}

\setlength{\headheight}{10pt}

\def \TITLE{GRADUATION THESIS}
\def \AUTHOR{Trần Văn A}

% ===================================================
\titleformat{\chapter}[hang]{\centering\bfseries}{CHAPTER \thechapter.\ }{0pt}{}[]

\titleformat
    {\chapter} % command
    [hang] % shape
    {\centering\bfseries} % format
    {CHAPTER \thechapter.\ } % label
    {0pt} %sep
    {} % before
    [] % after
\titlespacing*{\chapter}{0pt}{-20pt}{20pt}

\titleformat
    {\section} % command
    [hang] % shape
    {\bfseries} % format
    {\thechapter.\arabic{section}\ \ \ \ } % label
    {0pt} %sep
    {} % before
    [] % after
\titlespacing{\section}{0pt}{\parskip}{0.5\parskip}

\titleformat
    {\subsection} % command
    [hang] % shape
    {\bfseries} % format
    {\thechapter.\arabic{section}.\arabic{subsection}\ \ \ \ } % label
    {0pt} %sep
    {} % before
    [] % after
\titlespacing{\subsection}{30pt}{\parskip}{0.5\parskip}

\renewcommand\thesubsubsection{\alph{subsubsection}}
\titleformat
    {\subsubsection} % command
    [hang] % shape
    {\bfseries} % format
    {\alph{subsubsection}, \ } % label
    {0pt} %sep
    {} % before
    [] % after
\titlespacing{\subsubsection}{50pt}{\parskip}{0.5\parskip}

% \newcommand{\titlesize}{\fontsize{18pt}{23pt}\selectfont}
% \newcommand{\subtitlesize}{\fontsize{16pt}{21pt}\selectfont}
% \titleclass{\part}{top}
% \titleformat{\part}[display]
%   {\normalfont\huge\bfseries}{\centering}{20pt}{\Huge\centering}
% \titlespacing{\part}{0pt}{em}{1em}
% \titlespacing{\section}{0pt}{\parskip}{0.5\parskip}
% \titlespacing{\subsection}{0pt}{\parskip}{0.5\parskip}
% \titlespacing{\subsubsection}{0pt}{\parskip}{0.5\parskip}



% ===================================================
\usepackage{hyperref}
\hypersetup{pdfborder = {0 0 0}}
\hypersetup{pdftitle={\TITLE},
	pdfauthor={\AUTHOR}}

\usepackage[all]{hypcap} % Cho phép tham chiếu chính xác đến hình ảnh và bảng biểu

\graphicspath{{figures/}{../figures/}} % Thư mục chứa các hình ảnh

\counterwithin{figure}{chapter} % Đánh số hình ảnh kèm theo chapter. Ví dụ: Hình 1.1, 1.2,..

\title{\bf \TITLE}
\author{\AUTHOR}

\setcounter{secnumdepth}{3} % Cho phép subsubsection trong report
% \setcounter{tocdepth}{3} % Chèn subsubsection vào bảng mục lục

\theoremstyle{definition}
\newtheorem{example}{Example}[chapter] % Định nghĩa môi trường ví dụ

\onehalfspacing
\setlength{\parskip}{6pt}
\setlength{\parindent}{15pt}



% =========================== BODY ===============
\begin{document}
% \newgeometry{top=2cm, bottom=2cm, left=2cm, right=2cm}
% \subfile{Cover} % Phần bìa
\subfile{Cover} % Phần bìa lót
% \restoregeometry

% ===================================================
\pagestyle{empty} % Header và footer rỗng

\newpage
\pagenumbering{gobble} % Xóa page numbering ở cuối trang
\subfile{Chapter/0_2_Acknowledgement}

\newpage
\pagenumbering{gobble} % Xóa page numbering ở cuối trang
\subfile{Chapter/0_3_Abstract.tex}


% ===================================================
\newpage
\pagenumbering{gobble} % Xóa page numbering ở cuối trang
\renewcommand*\contentsname{TABLE OF CONTENTS}
\titlecontents{chapter}
    [0.0cm]             % left margin
    {\bfseries\vspace{0.3cm}}                  % above code
    {{\bfseries{\scshape} CHAPTER \thecontentslabel.\ }} % numbered format
    {}         % unnumbered format
    {\titlerule*[0.3pc]{.}\contentspage}         % filler-page-format, e.g dots

\titlecontents{section}
    [0.0cm]             % left margin
    {\vspace{0.3cm}}                  % above code
    {\thecontentslabel \ } % numbered format
    {}         % unnumbered format
    {\titlerule*[0.3pc]{.}\contentspage}         % filler-page-format, e.g dots

\titlecontents{subsection}
    [1.0cm]             % left margin
    {\vspace{0.3cm}}                  % above code
    {\thecontentslabel \ } % numbered format
    {}         % unnumbered format
    {\titlerule*[0.3pc]{.}\contentspage}         % filler-page-format, e.g dots

% ---- TABLE OF CONTENTS ----
\addtocontents{toc}{\protect\thispagestyle{empty}}
\tableofcontents
\thispagestyle{empty}
\clearpage

\pagenumbering{roman}

% ---- LIST OF FIGURES ----
\renewcommand{\listfigurename}{LIST OF FIGURES}
{\let\oldnumberline\numberline
 \renewcommand{\numberline}{Figure~\oldnumberline}
 \listoffigures}
\clearpage

% ---- LIST OF TABLES ----
\renewcommand{\listtablename}{LIST OF TABLES}
{\let\oldnumberline\numberline
 \renewcommand{\numberline}{Table~\oldnumberline}
 \listoftables}
\clearpage

% ---- LIST OF ABBREVIATIONS ----
\renewcommand*{\acronymname}{LIST OF ABBREVIATIONS}
\renewcommand*{\entryname}{Abbreviation}
\renewcommand*{\descriptionname}{Definition}
\glsaddall
\printnoidxglossary[type=\acronymtype, title={LIST OF ABBREVIATIONS}, nonumberlist]
\clearpage

\renewcommand\appendixname{APPENDIX}
\renewcommand\appendixpagename{APPENDIX}
\renewcommand\appendixtocname{APPENDIX}
\renewcommand{\figurename}{Figure}
\renewcommand{\tablename}{Table}
\renewcommand{\chaptername}{CHAPTER}

% ===================================================
\pagenumbering{arabic}
\pagestyle{fancy}
\fancyhf{}
\fancyhead[RE, LO]{\leftmark}
%\fancyhead[LE]{\rightmark}
\fancyfoot[RE, LO]{\thepage}

\chapter{INTRODUCTION}
\subfile{Chapter/1_Introduction} % Phần mở đầu

\newpage
\chapter{BACKGROUND AND RELATED WORK}
\subfile{Chapter/2_Literature_review}

\newpage
\chapter{PROPOSED SOLUTION}
\subfile{Chapter/3_Methodology}

\newpage
\chapter{EVALUATION AND RESULTS}
\subfile{Chapter/5_Numerical_results}

\newpage
\chapter{CONCLUSIONS AND FUTURE WORK}
\subfile{Chapter/6_Conclusions}

% ===================================================
\newpage
\renewcommand\bibname{REFERENCE}
\printbibliography
\phantomsection\addcontentsline{toc}{chapter}{REFERENCE}

\end{document}
